% !TEX root = ../main.tex

\chapter{Hasil Pengujian}

% ----------------------------------------------------------------
\section{Pengaruh \textit{Depth} dan \textit{Width}}

\subsection{Variasi \textit{Width} (Depth Tetap)}

\lipsum[1]

\begin{table}[H]
    \centering
    \caption{Hasil Prediksi Variasi Width}
    \begin{tabular}{llll}
        \toprule
        \textbf{Konfigurasi} & \textbf{Arsitektur} & \textbf{Train Loss} & \textbf{Val Loss} \\
        \midrule
        Width kecil  & [28, 16, 8, 1] & 0.4820 & 0.5056 \\
        Width sedang & [28, 32, 16, 1] & 0.4735 & 0.5124  \\
        Width besar  & [28, 64, 32, 1] & 0.4644 & 0.5091 \\
        \bottomrule
    \end{tabular}
    \label{tab:width}
\end{table}

\begin{figure}[H]
    \centering
    \includegraphics[width=0.4\textwidth]{ad_width_loss.png}
    \includegraphics[width=0.4\textwidth]{ad_width_val.png}
    \caption{Grafik training dan validation loss variasi width}
    \label{fig:width-loss}
\end{figure}

\subsection{Variasi \textit{Depth} (Width Tetap)}

\lipsum[1]

\begin{table}[H]
    \centering
    \caption{Hasil Prediksi Variasi Depth}
    \begin{tabular}{llll}
        \toprule
        \textbf{Konfigurasi} & \textbf{Arsitektur} & \textbf{Train Loss} & \textbf{Val Loss} \\
        \midrule
        Shallow & [28, 32, 1] & 0.4787 & 0.5062 \\
        Medium  & [28, 32, 32, 1] & 0.4684 & 0.5160 \\
        Deep    & [28, 32, 32, 32, 1] & 0.4644 & 0.5162 \\
        \bottomrule
    \end{tabular}
    \label{tab:depth}
\end{table}

\begin{figure}[H]
    \centering
    \includegraphics[width=0.4\textwidth]{ad_depth_loss.png}
    \includegraphics[width=0.4\textwidth]{ad_depth_val.png}
    \caption{Grafik training dan validation loss variasi depth}
    \label{fig:depth-loss}
\end{figure}


% ----------------------------------------------------------------
\section{Pengaruh Fungsi Aktivasi}

\lipsum[1]

\begin{table}[H]
    \centering
    \caption{Hasil Prediksi Variasi Fungsi Aktivasi}
    \begin{tabular}{lll}
        \toprule
        \textbf{Fungsi Aktivasi} & \textbf{Train Loss} & \textbf{Val Loss} \\
        \midrule
        Linear  & 0.4818 & 0.5049 \\
        ReLU    & 0.4757 & 0.5105 \\
        Sigmoid & 0.4873 & 0.5015 \\
        Tanh    & 0.4805 & 0.5051 \\
        \bottomrule
    \end{tabular}
    \label{tab:activation-exp}
\end{table}

\begin{figure}[H]
    \centering
    \includegraphics[width=0.4\textwidth]{ad_act_loss.png}
    \includegraphics[width=0.4\textwidth]{ad_act_val.png}
    \caption{Grafik training dan validation loss variasi fungsi aktivasi}
    \label{fig:activation-loss}
\end{figure}

\begin{figure}[H]
    \centering
    \includegraphics[width=0.8\textwidth]{ad_act_weight_lin.png}
    \caption{Grafik distribusi bobot variasi fungsi aktivasi linear}
    \label{fig:activation-weight-dist-lin}
\end{figure}

\begin{figure}[H]
    \centering
    \includegraphics[width=0.8\textwidth]{ad_act_weight_relu.png}
    \caption{Grafik distribusi bobot variasi fungsi aktivasi relu}
    \label{fig:activation-weight-dist-relu}
\end{figure}

\begin{figure}[H]
    \centering
    \includegraphics[width=0.8\textwidth]{ad_act_weight_sig.png}
    \caption{Grafik distribusi bobot variasi fungsi aktivasi sigmoid}
    \label{fig:activation-weight-dist-sig}
\end{figure}

\begin{figure}[H]
    \centering
    \includegraphics[width=0.8\textwidth]{ad_act_weight_tanh.png}
    \caption{Grafik distribusi bobot variasi fungsi aktivasi tanh}
    \label{fig:activation-weight-dist-tanh}
\end{figure}

\begin{figure}[H]
    \centering
    \includegraphics[width=0.8\textwidth]{ad_act_grad_lin.png}
    \caption{Grafik distribusi gradien variasi fungsi aktivasi linear}
    \label{fig:activation-grad-dist-lin}
\end{figure}

\begin{figure}[H]
    \centering
    \includegraphics[width=0.8\textwidth]{ad_act_grad_relu.png}
    \caption{Grafik distribusi gradien variasi fungsi aktivasi relu}
    \label{fig:activation-grad-dist-relu}
\end{figure}

\begin{figure}[H]
    \centering
    \includegraphics[width=0.8\textwidth]{ad_act_grad_sig.png}
    \caption{Grafik distribusi gradien variasi fungsi aktivasi sigmoid}
    \label{fig:activation-grad-dist-sig}
\end{figure}

\begin{figure}[H]
    \centering
    \includegraphics[width=0.8\textwidth]{ad_act_grad_tanh.png}
    \caption{Grafik distribusi gradien variasi fungsi aktivasi tanh}
    \label{fig:activation-grad-dist-tanh}
\end{figure}

% ----------------------------------------------------------------
\section{Pengaruh \textit{Learning Rate}}

\lipsum[1]

\begin{table}[H]
    \centering
    \caption{Hasil Prediksi Variasi Learning Rate}
    \begin{tabular}{lll}
        \toprule
        \textbf{Learning Rate} & \textbf{Train Loss} & \textbf{Val Loss} \\
        \midrule
        $\eta = 0.0001$ & 0.4930 & 0.5126 \\
        $\eta = 0.001$  & 0.4725 & 0.5144 \\
        $\eta = 0.01$   & 0.4631 & 0.5251 \\
        \bottomrule
    \end{tabular}
    \label{tab:lr}
\end{table}

\begin{figure}[H]
    \centering
    \includegraphics[width=0.4\textwidth]{ad_lr_loss.png}
    \includegraphics[width=0.4\textwidth]{ad_lr_val.png}
    \caption{Grafik training dan validation loss variasi learning rate}
    \label{fig:lr-loss}
\end{figure}

\begin{figure}[H]
    \centering
    \includegraphics[width=0.8\textwidth]{ad_lr_weight_small.png}
    \caption{Grafik distribusi bobot variasi learning rate 0.0001}
    \label{fig:lr-weight-dist-small}
\end{figure}

\begin{figure}[H]
    \centering
    \includegraphics[width=0.8\textwidth]{ad_lr_weight_med.png}
    \caption{Grafik distribusi bobot variasi learning rate 0.001}
    \label{fig:lr-weight-dist-med}
\end{figure}

\begin{figure}[H]
    \centering
    \includegraphics[width=0.8\textwidth]{ad_lr_weight_big.png}
    \caption{Grafik distribusi bobot variasi learning rate 0.01}
    \label{fig:lr-weight-dist-big}
\end{figure}

\begin{figure}[H]
    \centering
    \includegraphics[width=0.8\textwidth]{ad_lr_grad_small.png}
    \caption{Grafik distribusi gradien variasi learning rate 0.0001}
    \label{fig:lr-grad-dist-small}
\end{figure}

\begin{figure}[H]
    \centering
    \includegraphics[width=0.8\textwidth]{ad_lr_grad_med.png}
    \caption{Grafik distribusi gradien variasi learning rate 0.001}
    \label{fig:lr-grad-dist-med}
\end{figure}

\begin{figure}[H]
    \centering
    \includegraphics[width=0.8\textwidth]{ad_lr_grad_big.png}
    \caption{Grafik distribusi gradien variasi learning rate 0.01}
    \label{fig:lr-grad-dist-big}
\end{figure}

% ----------------------------------------------------------------
\section{Pengaruh Regularisasi}

\lipsum[1]

\begin{table}[H]
    \centering
    \caption{Hasil Prediksi Variasi Regularisasi}
    \begin{tabular}{lll}
        \toprule
        \textbf{Regularisasi} & \textbf{Train Loss} & \textbf{Val Loss} \\
        \midrule
        Tanpa regularisasi & 0.4742 & 0.5152 \\
        L1 & 0.5267 & 0.5005 \\
        L2 & 0.4932 & 0.5070 \\
        \bottomrule
    \end{tabular}
    \label{tab:reg}
\end{table}

\begin{figure}[H]
    \centering
    \includegraphics[width=0.4\textwidth]{ad_reg_loss.png}
    \includegraphics[width=0.4\textwidth]{ad_reg_val.png}
    \caption{Grafik training dan validation loss variasi regularisasi}
    \label{fig:reg-loss}
\end{figure}

\begin{figure}[H]
    \centering
    \includegraphics[width=0.8\textwidth]{ad_weight_noreg.png}
    \caption{Grafik distribusi bobot variasi tanpa regularisasi}
    \label{fig:reg-weight-dist-noreg}
\end{figure}

\begin{figure}[H]
    \centering
    \includegraphics[width=0.8\textwidth]{ad_weight_l1.png}
    \caption{Grafik distribusi bobot variasi regularisasi l1}
    \label{fig:reg-weight-dist-l1}
\end{figure}

\begin{figure}[H]
    \centering
    \includegraphics[width=0.8\textwidth]{ad_weight_l2.png}
    \caption{Grafik distribusi bobot variasi regularisasi l2}
    \label{fig:reg-weight-dist-l2}
\end{figure}

\begin{figure}[H]
    \centering
    \includegraphics[width=0.8\textwidth]{ad_grad_noreg.png}
    \caption{Grafik distribusi gradien variasi tanpa regularisasi}
    \label{fig:reg-grad-dist-noreg}
\end{figure}

\begin{figure}[H]
    \centering
    \includegraphics[width=0.8\textwidth]{ad_grad_l1.png}
    \caption{Grafik distribusi gradien variasi regularisasi l1}
    \label{fig:reg-grad-dist-l1}
\end{figure}

\begin{figure}[H]
    \centering
    \includegraphics[width=0.8\textwidth]{ad_grad_l2.png}
    \caption{Grafik distribusi gradien variasi regularisasi l2}
    \label{fig:reg-grad-dist-l2}
\end{figure}

% ----------------------------------------------------------------
\section{Perbandingan dengan \textit{sklearn} MLP}

\lipsum[1]

\begin{table}[H]
    \centering
    \caption{Perbandingan Hasil Prediksi}
    \begin{tabular}{lllll}
        \toprule
        \textbf{Model} & \textbf{Akurasi} & \textbf{Presisi} & \textbf{Recall} & \textbf{F1-Score} \\
        \midrule
        Implementasi FFNN & -- & -- & -- & -- \\
        Implementasi Autodiff & 0.7495 & 0.7906 & 0.8066 & 0.7985 \\
        sklearn \texttt{MLPClassifier} & 0.7540 & 0.7759 & 0.8440 & 0.8085 \\
        \bottomrule
    \end{tabular}
    \label{tab:sklearn}
\end{table}